\begin{proposition}[Near-bridge separated energy]
\label{prop:near-bridge-energy}
Let \(X\) be sufficiently large, put
\[
 H=\lfloor\sqrt{2X}\rfloor,\qquad 2\le K\le H,
\]
and, for every prime \(X<p\le2X\), let \(\pi_p\) be the projective
Ap\'ery orbit on \(\{0,\ldots,p-1\}\).  For each projective state, select
any reflection-invariant collection of four consecutive occurrences whose
span is at most \(\lfloor\sqrt p\rfloor\), and quotient the selected windows
by reflection.  Call two quotient orbits separated if the closed interval
of every oriented representative of one is disjoint from the closed interval
of every oriented representative of the other.

For an ordered separated pair of distinct quotient orbits, choose
deterministically an oriented representative
\[
 \gamma=(x,z_1,z_2,x+s)
\]
of the first orbit and an oriented representative
\[
 \delta=(y,w_1,w_2,y+t)
\]
of the second orbit such that \(x+s<y\), minimizing first
\(G=y-x-s\) and then the pair \((x,y)\) lexicographically.  Let
\(E_p^{\mathrm{near}}(K)\) count the ordered separated orbit pairs for
which this selected bridge satisfies \(G\le K\).  Then
\[
 \sum_{X<p\le2X}E_p^{\mathrm{near}}(K)\ll H^2K^2.
\]
In particular, for
\[
 K=\left\lfloor\frac{H}{\sqrt{\log X}}\right\rfloor
\]
one has
\[
 \sum_{X<p\le2X}E_p^{\mathrm{near}}(K)
 \ll\frac{X^2}{\log X}.
\]
\end{proposition}

\begin{proof}
Write \(R(I)=[p-1-v,p-1-u]\) for the reflection of an interval
\(I=[u,v]\).  The four cross-orbit intervals attached to a separated
orbit pair are pairwise disjoint across the two orbits.  There is always
a representative of the first orbit lying strictly to the left of a
representative of the second.  Indeed, if an interval of the second lies
to the left of an interval of the first, applying \(R\) reverses their
order and supplies the required representatives.  The deterministic
choice in the statement is therefore well-defined.

The map from an ordered orbit pair to
\begin{equation}
 (p,x,s,G)
 \label{eq:near-bridge-tuple}
\end{equation}
is injective.  To see this, put \(q=\pi_p(x)\) and
\(y=x+s+G\).  Since \(\gamma\) consists of four consecutive occurrences
of \(q\), its two internal indices are the first two occurrences of \(q\)
after \(x\), and hence are uniquely determined by \(p\) and \(x\).
The endpoint \(x+s\) verifies which third return is selected.  The same
argument at \(y\) uniquely determines \(\delta\), including its span
\(t\).  Thus the second span need not be retained as a certificate
parameter.  This use of consecutiveness is essential; the assertion would
be false for arbitrary four-element subsets of a fiber.

Retain the positional-gap continuants
\[
 \begin{aligned}
 N_1(u)&=1,\qquad N_2(u)=P(u+1),\\
 N_{h+1}(u)&=P(u+h)N_h(u)-(u+h)^6N_{h-1}(u).
 \end{aligned}
\]
By Lemma~\ref{lem:gap-poly}, the exact return criterion on a
nonwrapping interval is
\[
 \pi_p(u)=\pi_p(u+h)\quad\Longleftrightarrow\quad
 N_h(u)\equiv0\pmod p.
\]
Consequently the tuple in \eqref{eq:near-bridge-tuple} supplies a common
root \(x\in\mathbf F_p\) of
\[
 N_s(x),\qquad N_G(x+s).
\]
The second chain also gives \(N_t(x+s+G)=0\), but dropping this third
equation only enlarges the set being counted.

For \(s,G\ge2\), put
\[
 \mathcal R_{s,G}
 =\operatorname{Res}_x\bigl(N_s(x),N_G(x+s)\bigr)\in\mathbf Z.
\]
Proposition~\ref{prop:meso-root-strip} states that every complex root of
\(N_s(x)\) has real part in \((-s,-1)\), whereas every complex root of
\(N_G(x+s)\) has real part in \((-s-G,-s-1)\).  The strips are disjoint,
so
\[
 \mathcal R_{s,G}\ne0.
\]

Let \(\nu_p(s,G)\) be the number of common roots of these two polynomials
in \(\mathbf F_p\).  Since \(s,G\le H<p\),
Lemma~\ref{lem:nonvanish} ensures that neither reduced polynomial is zero.
The reduction modulo \(p\) of their Sylvester matrix has nullity at least
\(\nu_p(s,G)\): the reduced greatest common divisor has degree at least
the number of distinct common roots.  Smith normal form over \(\mathbf Z\)
therefore gives
\begin{equation}
 \nu_p(s,G)\le v_p(\mathcal R_{s,G}).
 \label{eq:near-bridge-smith}
\end{equation}
This remains valid if a leading coefficient vanishes modulo \(p\), since
degree loss can only increase the Sylvester nullity.

It remains to bound the height.  The polynomial \(N_h\) is the determinant
of an \((h-1)\)-by-\((h-1)\) tridiagonal matrix whose nonzero entries have
coefficient \(\ell^1\)-norm \(O(h^3)\).  The Leibniz expansion therefore
gives
\[
 \log\lVert N_h\rVert_2\ll h\log(2h).
\]
Translation satisfies
\[
 \lVert f(x+s)\rVert_1
 \le \lVert f\rVert_1(1+s)^{\deg f}.
\]
Since \(\deg N_h=3(h-1)\), the standard resultant norm inequality yields,
uniformly for \(s,G\le H\),
\begin{equation}
 \log|\mathcal R_{s,G}|
 \ll sG\log(2H).
 \label{eq:near-bridge-height}
\end{equation}

Injectivity, \eqref{eq:near-bridge-smith}, and
\eqref{eq:near-bridge-height} now give
\[
 \begin{aligned}
 \sum_{X<p\le2X}E_p^{\mathrm{near}}(K)
 &\le
 \sum_{s\le H}\sum_{G\le K}\sum_{X<p\le2X}
 v_p(\mathcal R_{s,G})\\
 &\le
 \frac1{\log X}
 \sum_{s\le H}\sum_{G\le K}\log|\mathcal R_{s,G}|\\
 &\ll
 \frac{\log(2H)}{\log X}
 \left(\sum_{s\le H}s\right)
 \left(\sum_{G\le K}G\right)
 \ll H^2K^2.
 \end{aligned}
\]
Here pairs with \(s=1\) or \(G=1\) do not occur because \(N_1=1\).
The final assertion follows from \(H^2\le2X\).
\end{proof}
